% =============================================================================
% Chapter 10: Monitoring and Checkpointing
% =============================================================================
\chapter{Monitoring and Checkpointing}
\label{ch:monitoring-checkpointing}

\begin{objectives}
    \item Use Weights \& Biases for experiment tracking
    \item Implement robust checkpointing
    \item Resume training from checkpoints
    \item Ensure reproducibility with seeds
\end{objectives}

% -----------------------------------------------------------------------------
\section{Experiment Tracking with W\&B}
% -----------------------------------------------------------------------------

\textbf{Weights \& Biases} (W\&B) provides cloud-based experiment tracking.

\subsection{What W\&B Logs}

\TryLM logs:
\begin{itemize}
    \item \textbf{Metrics}: Loss, perplexity, gradient norm, learning rate
    \item \textbf{Config}: All hyperparameters
    \item \textbf{Samples}: Generated text examples during training
    \item \textbf{System}: GPU memory, utilization
\end{itemize}

\subsection{Enabling W\&B}

\begin{lstlisting}[style=shell]
# Login (one-time)
uv run wandb login

# Train with logging
uv run trylm train --wandb-project trylm --wandb-run my-experiment
\end{lstlisting}

\subsection{Implementation}

\begin{lstlisting}[style=python, caption={W\&B logging from trainer.py}]
import wandb

# Initialize
wandb.init(
    project=config.wandb_project,
    name=config.wandb_run_name,
    config=asdict(config),
)

# Log metrics
wandb.log({
    "train/loss": loss,
    "train/perplexity": math.exp(loss),
    "train/learning_rate": lr,
    "train/grad_norm": grad_norm,
}, step=step)

# Log generated samples
samples = generate_samples(model, tokenizer)
wandb.log({"samples": wandb.Table(data=samples)}, step=step)
\end{lstlisting}

% -----------------------------------------------------------------------------
\section{Checkpointing}
% -----------------------------------------------------------------------------

Checkpoints save training state so you can resume or use the model later.

\subsection{Checkpoint Contents}

\begin{lstlisting}[style=shell, numbers=none]
checkpoints/step_5000/
|-- model.safetensors      # Model weights
|-- training_state.pt      # Optimizer, scheduler, step
`-- config.yaml            # Full configuration
\end{lstlisting}

\subsection{Saving Checkpoints}

\begin{lstlisting}[style=python, caption={Checkpoint saving from trainer.py}]
from safetensors.torch import save_file

def save_checkpoint(self, path: Path, step: int) -> None:
    path.mkdir(parents=True, exist_ok=True)

    # Save model weights with safetensors
    state_dict = self.model.state_dict()
    save_file(state_dict, path / "model.safetensors")

    # Save training state
    torch.save({
        "step": step,
        "optimizer": self.optimizer.state_dict(),
        "scheduler": self.scheduler.state_dict() if self.scheduler else None,
        "scaler": self.scaler.state_dict() if self.scaler else None,
        "best_val_loss": self.best_val_loss,
    }, path / "training_state.pt")

    # Save config for reproducibility
    self.full_config.to_yaml(path / "config.yaml")
\end{lstlisting}

\subsection{Loading Checkpoints}

\begin{lstlisting}[style=python, caption={Checkpoint loading from trainer.py}]
from safetensors.torch import load_file

def load_checkpoint(self, path: Path) -> int:
    # Load model weights
    state_dict = load_file(path / "model.safetensors")
    self.model.load_state_dict(state_dict)

    # Load training state
    training_state = torch.load(path / "training_state.pt")
    self.optimizer.load_state_dict(training_state["optimizer"])
    if self.scheduler and training_state["scheduler"]:
        self.scheduler.load_state_dict(training_state["scheduler"])
    if self.scaler and training_state["scaler"]:
        self.scaler.load_state_dict(training_state["scaler"])

    return training_state["step"]
\end{lstlisting}

\subsection{Checkpoint Strategy}

\TryLM saves checkpoints at three points:
\begin{enumerate}
    \item \textbf{Regular intervals}: Every \texttt{save\_interval} steps (default: 1000)
    \item \textbf{Best model}: When validation loss improves
    \item \textbf{Final}: At training completion
\end{enumerate}

% -----------------------------------------------------------------------------
\section{Resuming Training}
% -----------------------------------------------------------------------------

\begin{lstlisting}[style=shell]
# Resume from a specific checkpoint
uv run trylm train --resume checkpoints/step_5000

# Resume from best checkpoint
uv run trylm train --resume checkpoints/best
\end{lstlisting}

Resuming restores:
\begin{itemize}
    \item Model weights
    \item Optimizer state (momentum, adaptive learning rates)
    \item Training step counter
    \item Learning rate scheduler state
    \item Mixed precision scaler state
\end{itemize}

% -----------------------------------------------------------------------------
\section{Reproducibility}
% -----------------------------------------------------------------------------

\subsection{Setting Seeds}

\begin{lstlisting}[style=python]
def set_seed(seed: int) -> None:
    random.seed(seed)
    np.random.seed(seed)
    torch.manual_seed(seed)
    torch.cuda.manual_seed_all(seed)
\end{lstlisting}

\subsection{Deterministic Operations}

\begin{lstlisting}[style=python]
# For full reproducibility (slower)
torch.backends.cudnn.deterministic = True
torch.backends.cudnn.benchmark = False
\end{lstlisting}

\begin{keypoint}
Even with seeds set, GPU operations may have non-deterministic behavior. For true reproducibility, save checkpoints frequently and validate on held-out data.
\end{keypoint}

% -----------------------------------------------------------------------------
\section{Summary}
% -----------------------------------------------------------------------------

\begin{summary}
    \item \textbf{W\&B} provides cloud experiment tracking with metrics, configs, and samples
    \item \textbf{Checkpoints} save model, optimizer, and scheduler state
    \item \textbf{Safetensors} format is safe and efficient for model weights
    \item \textbf{Resume training} restores complete state from checkpoint
    \item \textbf{Seeds} enable approximate reproducibility
\end{summary}
